\chapter*{Ringraziamenti}

Come nei migliori film di Christopher Nolan, siamo giunti al momento finale, ovvero al momento che qualsiasi studente universitario aspetta più di ogni altra cosa: la laurea magistrale. Un traguardo che non solo sancisce la fine di un percorso di studi, ma che, soprattutto, apre le porte a una nuova fase della vita e alla carriera lavorativa. Personalmente, sto vivendo questo ultimo periodo con un po' di rammarico. Nonostante il ``veleno buttato'' per gli esami in questi anni, credo che la vita dinamica, le sfide e le esperienze vissute durante questo percorso mi abbiano cambiato profondamente come persona. Mi rendo conto che, più degli esami superati e del Politecnico, saranno le persone incontrate, i luoghi vissuti e tutti quei momenti apparentemente insignificanti a rimanere impressi nella mia memoria.

Adesso però, bando alle ciance. Se questo capitolo si chiama ``Ringraziamenti'', la pianto di parlare di me e passiamo all'unica sezione di questa tesi che davvero interessa a qualcuno, in primis a me più di tutti.

Grazie alla mia famiglia. Mi avete visto scappare a mille chilometri da casa per inseguire i miei sogni e avete avuto fiducia in un diciottenne che era convinto della propria scelta. Vi confesso che, alla base di quella scelta, c'erano motivazioni molto vaghe e probabilmente non avevo nemmeno ben chiaro dove stessi andando. Però vi assicuro che tutti i sacrifici che avete fatto per sostenermi sono stati ripagati ampiamente. Mi avete dato la libertà di sbagliare, di mettermi alla prova e, soprattutto, di credere in me anche quando forse ero io il primo a non farlo. Se oggi sono arrivato fin qui, una parte enorme del merito è vostra.

Grazie a tutte le persone che hanno fatto parte di questo percorso, dagli amici del Sud a quelli del Nord. Siete così tanti che sarebbe impossibile citarvi tutti senza trasformare questi ringraziamenti in un'altra tesi. Siete stati il sale che ha dato sapore a questi anni: le serate improvvisate, le risate, i pianti, le scampagnate, i viaggi, le discussioni e tutte quelle piccole cose che hanno reso questo periodo della mia vita davvero speciale.

Rompo per un attimo la monotonia del ``grazie a'' per riflettere su quanto sia stato importante l'amore in ogni momento della mia vita. L'amore è assurdo, perché è impossibile da razionalizzare. Arriva, ti rapisce e riesce a rendere speciale anche la quotidianità più banale. Credo sia una delle poche cose che non abbiamo bisogno di capire fino in fondo per riconoscerne il valore. Per questo, invito chiunque legga o ascolti queste parole a fermarsi, almeno per un momento, e a riflettere su quanto l'amore possa essere la risposta a molte delle cose negative che possono attraversare le nostre vite. Io oggi mi fermo e ringrazio chi, ogni giorno, riesce a trasmettermi amore attraverso un semplice sorriso, uno sguardo, una parola o un semplice gesto. Grazie, Ale.

Grazie Torino. Non ti conoscevo quando ti ho vista per la prima volta, eppure è stato amore a prima vista. Mi hai accolto quando avevo ancora tutto da scoprire e, senza che me ne rendessi conto, sei diventata casa. Mi hai dato un sacco di cose preziose: il Collegio Einaudi, il gruppo scout Torino 110, i terroni e un'infinità di ricordi che porterò con me.

Grazie anche a te, Costa Azzurra. Anche se solo per un anno, mi hai dato la possibilità di rimettermi in gioco, di uscire ancora una volta dalla mia zona di comfort e di imparare una nuova lingua, il francese, rafforzando al tempo stesso quelle che già conoscevo, l'inglese e l'italiano. Mi mancheranno il mare dopo il lavoro, le serate a Nizza, il casinò e, soprattutto, tutte le persone stupende che ho avuto la fortuna di conoscere.

Dulcis in fundo, grazie a te, Gravina. Nonostante tutte le lamentele che avrei da farti, sei il posto da cui tutto è iniziato. Mi hai dato una famiglia, degli amici, le prime esperienze, il gruppo scout Gravina 1, le feste, i pranzi della domenica, i paesaggi suggestivi e quel senso di appartenenza che, quando sei lontano, capisci davvero quanto sia importante. E forse è proprio questa la cosa più preziosa che mi hai lasciato: la consapevolezza di avere sempre un posto da chiamare casa, indipendentemente da quanti chilometri mi separano da te. Probabilmente continuerò a lamentarmi di te, come ho sempre fatto, ma so che, in fondo, una parte di me resterà sempre lì.

\begin{flushright}
{\itshape Giovanni Russo}
\end{flushright}
